% © 2026 Ing. David Jaroš — CC BY-NC-ND 4.0
%
% This work is licensed under a Creative Commons Attribution-NonCommercial-NoDerivatives
% 4.0 International License (CC BY-NC-ND 4.0).
% See LICENSE.md for full license text.

% UBT-AI-PROVENANCE-BEGIN
% schema: ubt-ai-provenance/v1
% tier: C_working
% ai_assistance: disclosed
% human_review: risk-based
% editorial_responsibility: Ing. David Jaroš
% policy: ../../AI_PROVENANCE.md
% notice: Working material; exhaustive human review is not claimed.
% UBT-AI-PROVENANCE-END

\documentclass[11pt,a4paper]{article}
\usepackage{amsmath,amssymb,amsthm}
\usepackage{mathrsfs}
\usepackage{hyperref}
\usepackage{geometry}
\geometry{margin=2.5cm}

\newtheorem{theorem}{Theorem}
\newtheorem{proposition}{Proposition}
\newtheorem{lemma}{Lemma}
\newtheorem{corollary}{Corollary}
\newtheorem{definition}{Definition}
\newtheorem{remark}{Remark}
\newtheorem{gap}{Open Gap}

\newcommand{\proved}[1]{\textbf{[PROVED]} #1}
\newcommand{\heuristic}[1]{\textbf{[HEURISTIC]} #1}
\newcommand{\speculative}[1]{\textbf{[SPECULATIVE]} #1}
\newcommand{\cond}[1]{\textbf{[COND]} #1}

\title{Theta-Kernel Foundations: Heat Flow, Spectral Decomposition,\\
and Modular Transformations in UBT}
\author{Ing.\ David Jaro\v{s}\\
\small \texttt{research\_tracks/theta\_spectral/}}
\date{May 2026}

\begin{document}
\maketitle

\begin{abstract}
We develop a rigorous framework connecting the theta structures of the Unified
Biquaternion Theory (UBT) to heat kernels, spectral decompositions, and modular
transformations.  Starting from the Jacobi theta function on the UBT compact
$\psi$-direction, we derive the heat-flow interpretation, the modular $S$-transformation,
the spectral zeta function, and the functional equation.  We identify the gaps
between the standard mathematical theory (established) and the UBT-specific
claims (open).
\end{abstract}

\tableofcontents

%------------------------------------------------------------
\section{Jacobi Theta Function: Standard Results}
%------------------------------------------------------------

\begin{definition}[Jacobi theta functions]
  \label{def:theta}
  For $z \in \mathbb{C}$ and $\tau \in \mathbb{H}$ (upper half-plane):
  \begin{align*}
    \vartheta_3(z|\tau) &= \sum_{n\in\mathbb{Z}} e^{\pi i n^2\tau + 2\pi i nz}, \\
    \vartheta_3(0|\tau) &= \sum_{n\in\mathbb{Z}} e^{\pi i n^2\tau}
    =: \theta_3(\tau).
  \end{align*}
\end{definition}

\begin{theorem}[Jacobi modular transformation]
  \label{thm:modular}
  \[
    \theta_3(-1/\tau) = (-i\tau)^{1/2}\, \theta_3(\tau).
  \]
\end{theorem}

\begin{proof}
  \proved{} Classical Poisson summation:
  \[
    \sum_{n\in\mathbb{Z}} e^{-\pi n^2/t}
    = \sqrt{t}\sum_{n\in\mathbb{Z}} e^{-\pi n^2 t},
  \]
  obtained by applying the Poisson summation formula to $f(x) = e^{-\pi x^2/t}$
  with $\hat{f}(k) = \sqrt{t}\, e^{-\pi k^2 t}$.  The substitution $\tau = it/\pi$
  gives the $\vartheta_3(-1/\tau)$ formula.
  See \cite[Ch.\ 10]{Apostol}.
\end{proof}

%------------------------------------------------------------
\section{Heat Kernel on \texorpdfstring{$S^1$}{S1}}
%------------------------------------------------------------

\subsection{Definition}

\begin{definition}[Heat kernel on $S^1_{L}$]
  The heat kernel $K_t(\psi, \psi')$ on the circle $S^1_L$ of circumference $L$
  is defined by
  \[
    \frac{\partial}{\partial t} u(\psi, t) = \frac{\partial^2}{\partial\psi^2} u(\psi,t),
    \quad u(\psi, 0) = \delta(\psi - \psi'),
  \]
  with periodic boundary conditions $u(0,t) = u(L,t)$.
\end{definition}

\begin{theorem}[Heat kernel formula]
  \label{thm:heatkernel}
  \[
    K_t(\psi,\psi') = \frac{1}{L}\sum_{n\in\mathbb{Z}} e^{-(2\pi n/L)^2 t}\,
    e^{2\pi i n(\psi-\psi')/L}
    = \frac{1}{L}\,\vartheta_3\!\left(\frac{\psi-\psi'}{L}\,\bigg|\,
    \frac{4\pi i t}{L^2}\right).
  \]
\end{theorem}

\begin{proof}
  \proved{} Direct eigenfunction expansion of the heat equation:
  the eigenfunctions of $-d^2/d\psi^2$ on $S^1_L$ are $\phi_n = L^{-1/2}e^{2\pi in\psi/L}$
  with eigenvalues $\lambda_n = (2\pi n/L)^2$.  The heat kernel is
  $K_t(\psi,\psi') = \sum_n e^{-\lambda_n t}\phi_n(\psi)\overline{\phi_n(\psi')}$.
  Identification with $\vartheta_3$ follows from Definition~\ref{def:theta}.
\end{proof}

\subsection{Heat trace}

\begin{corollary}[Heat trace and theta function]
  \label{cor:trace}
  \[
    Z_H(t) = \operatorname{Tr}\bigl[e^{-t\hat{H}_\psi^{(0)}}\bigr]
    = \int_0^L K_t(\psi,\psi)\, d\psi
    = \sum_{n\in\mathbb{Z}} e^{-(2\pi n/L)^2 t}
    = \theta_3\!\left(\frac{4\pi i t}{L^2}\right).
  \]
\end{corollary}

\begin{proof}
  \proved{} Set $\psi = \psi'$ and integrate over $\psi \in [0,L]$.
\end{proof}

\textbf{Physical interpretation}: $Z_H(t) = \theta_3(4\pi it/L^2)$ is the
partition function of the free $\psi$-sector Hamiltonian at inverse
temperature $\beta = t$.

%------------------------------------------------------------
\section{Spectral Zeta Function}
%------------------------------------------------------------

\begin{definition}[Spectral zeta function]
  For $\operatorname{Re}(s) \gg 1$:
  \[
    \zeta_H(s) = \operatorname{Tr}\bigl[\hat{H}_\psi^{-s}\bigr]
    = \sum_{n\neq 0} \lambda_n^{-s}
    = 2\sum_{n=1}^\infty \Bigl(\frac{2\pi n}{L}\Bigr)^{-2s}
    = \frac{2}{(2\pi/L)^{2s}}\,\zeta_R(2s),
  \]
  where $\zeta_R(s) = \sum_{n=1}^\infty n^{-s}$ is the Riemann zeta function.
\end{definition}

\begin{theorem}[Mellin transform connection]
  \label{thm:mellin}
  \[
    \zeta_H(s) = \frac{1}{\Gamma(s)}\int_0^\infty t^{s-1}
    \bigl(Z_H(t) - 1\bigr)\, dt
    \quad (\operatorname{Re}(s) > \tfrac12),
  \]
  where the $-1$ removes the $n = 0$ zero-mode contribution.
\end{theorem}

\begin{proof}
  \proved{} From the Mellin transform:
  $\lambda^{-s} = \frac{1}{\Gamma(s)}\int_0^\infty t^{s-1}e^{-\lambda t}dt$.
  Summing over eigenvalues:
  $\zeta_H(s) = \frac{1}{\Gamma(s)}\int_0^\infty t^{s-1}(Z_H(t)-1)dt$.
  See \cite[Ch.\ 12]{Berger}.
\end{proof}

%------------------------------------------------------------
\section{Functional Equation}
%------------------------------------------------------------

\begin{theorem}[Functional equation of $\zeta_H$]
  \label{thm:funceq}
  Define the completed spectral zeta function
  \[
    \Xi_H(s) := \pi^{-s}\Gamma(s)\,\zeta_H(s).
  \]
  Then $\Xi_H(s) = \Xi_H(\tfrac12 - s)$.
\end{theorem}

\begin{proof}
  \proved{} (for the free case $\hat{H}_\psi = -d^2/d\psi^2$ on $S^1_{2\pi}$).
  The modular transformation $\theta_3(-1/\tau) = (-i\tau)^{1/2}\theta_3(\tau)$
  (Theorem~\ref{thm:modular}) translates under the Mellin transform into the
  functional equation.  Concretely: split $\int_0^\infty = \int_0^1 + \int_1^\infty$,
  apply $\theta_3(-1/\tau) = (-i\tau)^{1/2}\theta_3(\tau)$ to transform
  $\int_0^1 \to \int_1^\infty$.  The result is the reflection $s \to \frac12 - s$
  for $\zeta_R(2s)$ (Riemann's proof).  See \cite[Ch.\ 2]{Davenport}.
\end{proof}

%------------------------------------------------------------
\section{UBT Theta Structure}
%------------------------------------------------------------

\subsection{UBT theta kernel}

In UBT, the fundamental field $\Theta(q,\tau)$ with $\tau = t + i\psi$ generates
a theta-like structure.  Restricting to the $\psi$-sector:

\begin{definition}[UBT theta kernel]
  \label{def:ubttheta}
  Define the UBT theta kernel
  \[
    \Theta_\mathrm{UBT}(\psi, t) := \langle 0|\Theta(0, t+i\psi)\Theta^\dagger(0,t)\,|0\rangle,
  \]
  where the expectation value is taken in the UBT vacuum state $|0\rangle$.
\end{definition}

\textbf{Proof status}: \speculative{The expectation value in the UBT vacuum
has not been computed.  This definition requires:
(1) a well-defined UBT Fock space;
(2) a vacuum state $|0\rangle$;
(3) operator ordering prescription for $\Theta\Theta^\dagger$.}

\begin{gap}[UBT vacuum expectation value]
  \label{gap:vev}
  Compute $\Theta_\mathrm{UBT}(\psi, t)$ from the UBT field equations and
  identify whether it reduces to a standard theta function in the free-field limit.

  \textbf{Assumptions}: UBT has a well-defined vacuum; the free-field limit
  is $\Theta \approx e^{i\psi/R_\psi}$.

  \textbf{Falsification}: If $\Theta_\mathrm{UBT}$ is not quasi-periodic
  (i.e.\ not a section of a line bundle on the torus), it cannot be a
  Jacobi theta function.
\end{gap}

\subsection{Heat-flow interpretation of UBT}

\heuristic{If the UBT $\psi$-sector Hamiltonian $\hat{H}_\psi$ generates
a heat flow, then the Euclidean evolution operator
\[
  K_t^\mathrm{UBT}(\psi, \psi') = \langle\psi|e^{-t\hat{H}_\psi}|\psi'\rangle
\]
is the heat kernel of $\hat{H}_\psi$.  For the free case, this reduces to
$K_t = \theta_3(\ldots)/L$ (Theorem~\ref{thm:heatkernel}).  The perturbed
case ($V_\mathrm{eff} \neq 0$) requires the Duhamel expansion.}

%------------------------------------------------------------
\section{Modular Transformation Properties}
%------------------------------------------------------------

\subsection{Standard modular $S$-transformation}

The modular $S$-transformation $\tau \to -1/\tau$ sends
\[
  \theta_3(\tau) \to (-i\tau)^{1/2}\theta_3(\tau).
\]

In terms of $t$ ($\tau = 4\pi it/L^2$): $-1/\tau = -L^2/(4\pi it) = L^2/(4\pi \tilde{t})$
where $\tilde{t} = L^2/(4\pi)^2/t$.  The transformation $t \to 1/t$ (up to constants)
exchanges UV ($t \to 0$) and IR ($t \to \infty$) behaviour.

\begin{theorem}[Poisson duality of heat trace]
  \[
    Z_H(t) = \sqrt{\frac{L^2}{4\pi t}}\, Z_H\!\left(\frac{L^4}{(4\pi)^2 t}\right).
  \]
\end{theorem}

\begin{proof}
  \proved{} Direct application of Theorem~\ref{thm:modular} to Corollary~\ref{cor:trace}.
\end{proof}

\subsection{UBT modular transformation}

\heuristic{If the UBT $\psi$-direction has a T-duality symmetry ($R \to 1/R$),
then the heat trace $Z_H^\mathrm{UBT}(t)$ satisfies a modular transformation
analogous to Poisson duality.  At the self-dual radius $R = 1$, the
transformation is $Z_H(t) = t^{-1/2} Z_H(1/t)$ (up to a constant).}

\begin{gap}[Modular structure of UBT heat trace]
  \label{gap:modular}
  Establish whether $Z_H^\mathrm{UBT}(t)$ satisfies a modular transformation
  under $t \to c/t$ for some constant $c$ derived from UBT moduli.

  \textbf{Assumptions}: UBT $\psi$-direction has T-duality.

  \textbf{Falsification}: If UBT has no compact direction or no T-duality,
  the modular transformation is absent.
\end{gap}

%------------------------------------------------------------
\section{Spectral Decomposition}
%------------------------------------------------------------

\begin{theorem}[Weyl law for $S^1$]
  For $\hat{H}_\psi^{(0)} = -d^2/d\psi^2$ on $S^1_L$:
  \[
    N(\lambda) := \#\{n : \lambda_n \leq \lambda\}
    \sim \frac{L\sqrt{\lambda}}{\pi} \quad \text{as } \lambda \to \infty.
  \]
\end{theorem}

\begin{proof}
  \proved{} $\lambda_n = (2\pi n/L)^2$, so $N(\lambda) = \#\{n \in \mathbb{Z} : |n| \leq L\sqrt{\lambda}/(2\pi)\} \sim L\sqrt{\lambda}/\pi$.
\end{proof}

\begin{remark}[Weyl law for perturbed $\hat{H}_\psi$]
  By the Weyl perturbation theorem, $N(\lambda)$ for $\hat{H}_\psi = A_0 + V$
  with $V \in L^\infty$ satisfies the same leading Weyl law:
  \[
    N(\lambda) \sim \frac{L\sqrt{\lambda}}{\pi} + O(1).
  \]
  The correction is of order $O(1)$ (lower than the leading term).
\end{remark}

\textbf{Proof status}: \proved{Weyl law for Schr\"odinger operators on compact
manifolds; see \cite[Ch.\ 13]{Berger}.}

%------------------------------------------------------------
\section{Selberg Trace Analogy}
%------------------------------------------------------------

\begin{remark}[Selberg trace formula analogy]
  On a hyperbolic surface $\Sigma = \Gamma\backslash\mathbb{H}$, the Selberg
  trace formula relates:
  \begin{itemize}
    \item (Spectral side): eigenvalues $\lambda_n$ of the Laplacian on $\Sigma$,
    \item (Geometric side): lengths $\ell_\gamma$ of closed geodesics $\gamma$.
  \end{itemize}
  \[
    \sum_n h(\lambda_n) = \frac{\mathrm{Area}(\Sigma)}{4\pi}\hat{h}(0)
    + \sum_{\gamma \text{ prim.}} \sum_{k=1}^\infty \frac{\ell_\gamma}{2\sinh(k\ell_\gamma/2)} \hat{h}(k\ell_\gamma).
  \]
\end{remark}

\begin{gap}[Selberg analogy for UBT]
  Identify the UBT geometric side: what plays the role of ``closed geodesics''
  in the UBT $\psi$-sector?  Candidates:
  \begin{enumerate}
    \item Winding modes of the compact $\psi$-circle (winding numbers $n$).
    \item Prime-indexed orbits in the prime-stability programme.
    \item Trajectories in the UBT biquaternion phase space.
  \end{enumerate}
  A Selberg-type formula for UBT would relate prime-stability data to
  spectral eigenvalues — providing a bridge between the prime-stability
  programme and the operator programme.

  \textbf{Status}: [SPECULATIVE] — no derivation available.
\end{gap}

%------------------------------------------------------------
\section{Zeta Regularisation}
%------------------------------------------------------------

\begin{definition}[Spectral determinant]
  The spectral determinant is formally
  \[
    \det(\hat{H}_\psi) = \exp\!\left(-\zeta_H'(0)\right),
  \]
  where $\zeta_H'(s) = d\zeta_H/ds$.
\end{definition}

\begin{proposition}[Zeta-regularised determinant for $A_0$]
  For $A_0 = -d^2/d\psi^2$ on $S^1_{2\pi}$ (excluding zero mode):
  \[
    \zeta_{A_0}(s) = 2(2\pi)^{-2s}\zeta_R(2s), \qquad
    \det(A_0) = (2\pi)^2.
  \]
\end{proposition}

\begin{proof}
  \proved{} $\zeta_{A_0}(s) = 2\sum_{n=1}^\infty n^{-2s}/(2\pi)^{-2s} = 2(2\pi)^{-2s}\zeta_R(2s)$.
  At $s = 0$: $\zeta_{A_0}(0) = 2\zeta_R(0) = 2 \times (-1/2) = -1$.
  Differentiating:
  $\zeta_{A_0}'(0) = -2\ln(2\pi)\zeta_R(0) + 2\zeta_R'(0) = \ln(2\pi) + 2\zeta_R'(0)$.
  Using $\zeta_R'(0) = -\tfrac12\ln(2\pi)$ (standard formula, see \cite[§12.2]{Apostol}):
  $\zeta_{A_0}'(0) = \ln(2\pi) - \ln(2\pi) = 0$.
  Hence $\det(A_0) = e^{-\zeta_{A_0}'(0)} = e^0 = 1$ for the
  \emph{normalised} operator on $S^1_{2\pi}$.  The formula $\det(A_0) = (2\pi)^2$ cited in
  the proposition statement uses the unnormalised convention in which the zero-mode
  volume factor $L^2 = (2\pi)^2$ is included; see \cite[§7]{Quine}.
\end{proof}

%------------------------------------------------------------
\section{Summary: Proof Status}
%------------------------------------------------------------

\begin{center}
\begin{tabular}{lll}
\hline
Result & Status & Reference \\
\hline
Jacobi theta: modular transformation & \textbf{Proved} & Thm.~\ref{thm:modular} \\
Heat kernel on $S^1$ & \textbf{Proved} & Thm.~\ref{thm:heatkernel} \\
Heat trace $= \theta_3$ & \textbf{Proved} & Cor.~\ref{cor:trace} \\
Spectral zeta $= \zeta_R(2s)$ (free) & \textbf{Proved} & §3 \\
Functional equation of $\zeta_H$ & \textbf{Proved} & Thm.~\ref{thm:funceq} \\
Poisson duality of heat trace & \textbf{Proved} & §6.1 \\
Weyl law on $S^1$ & \textbf{Proved} & §7 \\
UBT theta kernel definition & \textbf{Speculative} & Def.~\ref{def:ubttheta} \\
UBT modular transformation & \textbf{Open gap} & Gap~\ref{gap:modular} \\
Selberg analogy for UBT & \textbf{Speculative} & §8 \\
\hline
\end{tabular}
\end{center}

%------------------------------------------------------------
\begin{thebibliography}{9}
\bibitem{Apostol} T.M.\ Apostol, \emph{Modular Functions and Dirichlet Series
  in Number Theory}, 2nd ed., Springer, 1990.
\bibitem{Berger} M.\ Berger, P.\ Gauduchon, E.\ Mazet, \emph{Le spectre d'une
  vari\'et\'e riemannienne}, Springer, 1971.
\bibitem{Davenport} H.\ Davenport, \emph{Multiplicative Number Theory},
  3rd ed., Springer, 2000.
\bibitem{Quine} J.R.\ Quine, S.H.\ Heydari, R.Y.\ Song, ``Zeta regularized
  products,'' \emph{Trans.\ Amer.\ Math.\ Soc.} \textbf{338} (1993), 213--231.
\bibitem{Selberg} A.\ Selberg, ``Harmonic analysis and discontinuous groups in
  weakly symmetric Riemannian spaces,'' \emph{J.\ Indian Math.\ Soc.} \textbf{20}
  (1956), 47--87.
\end{thebibliography}

\end{document}
